\section{相关工作}

\paraf{RL 框架优化。}
LLM 强化学习的价值推动了大量 RL 框架优化研究~\citep{yao2023deepspeed,sheng2025hybridflow,hu2024openrlhf,zhong2025optimizing,zhong2025streamrl,fu2025areal,he2025history,wang2025distflow,lei2024puzzle,shen2024nemo,han2025asyncflow}。OpenRLHF~\citep{hu2024openrlhf} 和 verl~\citep{sheng2025hybridflow} 分别探索了解耦与共置架构。后续一些框架又提出融合~\citep{zhong2025optimizing}、流水化~\citep{zhong2025optimizing}、流式执行~\citep{zhong2025streamrl} 以及异步执行~\citep{fu2025areal} 等机制来提高系统效率。此外，还有一些优化工作提出了请求级调度、负载均衡和 speculative decoding~\citep{gao2025rollpacker,jinefficient}，以进一步加速 rollout 阶段。尽管这些工作很有价值，但现有框架通常仍把外部调用视为次要问题，要么忽略其资源管理，要么默认采用静态过量预留。因此，在轨迹更长、环境交互更复杂的智能体 RL 中，会造成明显资源浪费。

\parabf{智能体资源管理。}
随着智能体 RL 复杂度不断提高，对工具与奖励所依赖外部资源的管理也变得更加重要。Mars~\citep{zhu2026mars} 等近期工作已经开始关注奖励计算效率，并把奖励计算部署为独立服务。Mars 利用“请求进入、批量输出”的特征，把请求级延迟要求放宽为批级约束。尽管这对奖励计算有效，但 \sysname 提供了更通用的方案：它为所有类型的工具使用与奖励服务建立了统一的动作级平台，覆盖 API 调用、代码执行与模拟器等多种场景。MegaFlow~\citep{zhang2026megaflow}、SkyRL-Agent~\citep{cao2025skyrl} 和 ROCK~\citep{wang2025let} 则对智能体服务进行了抽象，并借鉴云计算方法来管理工具。通过把工具的可扩展性与共享模式纳入建模，\sysname 能对工具调用和奖励服务做更有针对性的优化。

\parabf{资源自动伸缩。}
自动伸缩是云计算的核心能力之一，通常采用响应式、预测式或混合式方法~\citep{lorido2014review,wu2023xron,wang2024autothrottle,wu2024can,zhang2024jolteon} 来满足服务等级目标（SLO）。在 LLM 场景中，相关系统~\citep{zhang2025blitzscale,yu2025lambda,xiang2025aegaeon,zhao2025seallm} 已经面向大规模推理工作负载优化了自动伸缩。针对 RL 场景，Mars~\citep{zhu2026mars} 又引入了基于历史的伸缩与超时感知机制，用于在 batch 边界上预测并调整资源需求。\sysname 在此基础上进一步提出了动作级弹性调度器，以应对智能体 RL 动作的突发性和异构性。通过对动作可扩展性的显式建模，\sysname 可以最小化总 ACT，从而加速 RL 训练过程。
